\section*{Complete author proof (original Section4, Theorem3)}
\begin{proof}
The proof is constructive and done by a two-step process. In the first step (which is fairly standard and works in many computational models) we move all the $\varepsilon$ with negative exponents to edges adjacent to the source. The second step then uses the monotonicity.
Given $\Gamma$ with affine linear labels over $\R[\varepsilon,\varepsilon^{-1}]_+$ we repeat the following process until all labels that contain an $\varepsilon$ with a negative exponent are incident to the source vertex.
\begin{itemize}
\item Let $e$ be an edge whose label contains $\varepsilon$ with a negative exponent $-i<0$. Moreover, assume that $e$ is not incident to the source vertex. Let $v$ be the start vertex of $e$. We rescale all edges outgoing of $v$ with $\varepsilon^i$ and we rescale all edges incoming to $v$ with $\varepsilon^{-i}$.
\end{itemize}
If we always choose the edge with the highest layer, then it is easy to see that this process terminates. Since every path from the source to the sink that goes through a vertex $v$ must use exactly one edge that goes into $v$ and exactly one edge that comes out of $v$, throughout the process the value of $\Gamma$ does not change. We finish this first phase by taking the highest negative power $i$ among all labels of edges that are incident to the source and then rescale all these edges with $\varepsilon^i$. The resulting ABP $\Gamma^i$ computes $\varepsilon^if_\epsilon$ and no label contains an $\varepsilon$ with negative exponent. We now start phase2 that transforms $\Gamma^i$ into $\Gamma^{i-1}$ that computes $\varepsilon^{i-1}f_\epsilon$ without introducing negative exponents of $\varepsilon$. We repeat phase2 until we reach $\Gamma^0$ in which we safely set $\varepsilon$ to0. Throughout the whole process we do not change the structure of the ABP and only rescale edge labels with powers of $\varepsilon$, which preserves monotonicity, so the proof is finished. It remains to show how $\Gamma^i$ can be transformed into $\Gamma^{i-1}$. An edge whose label is divisible by $\varepsilon$ is called an $\varepsilon$-edge. Consider the set $\Delta$ of vertices that are reachable from the source using only non $\varepsilon$-edges in $\Gamma^i$. The crucial insight is that since $\Gamma^i$ is monotone and computes a polynomial that is divisible by $\varepsilon$, we know that every path in $\Gamma^i$ from the source to the sink uses an $\varepsilon$-edge. Therefore $\Delta$ cannot contain the sink. We call a vertex in $\Delta$ whose outdegree is zero a leaf vertex. We repeat the following procedure until the source is the only leaf vertex:
\begin{itemize}
\item Let $v$ be a non-source leaf vertex in $\Delta$. We rescale all edges outgoing of $v$ with $\varepsilon^{-1}$ and we rescale all edges incoming to $v$ with $\varepsilon$.
\end{itemize}
It is easy to see that this process terminates with the source being the only leaf vertex. Since the source is a leaf vertex, all edges incident to the source are $\varepsilon$-edges. We divide all their labels by $\varepsilon$ to obtain $\Gamma^{i-1}$.
\end{proof}
